\begin{proposition}\label{prop:psef-cone-compact-base}
The pseudoeffective cone in $H^{1,1}_{\mathrm{BC}}(X,\mathbb R)$ is closed and convex and contains no nonzero line. For a fixed Gauduchon metric $\gamma$, the functional
\[
 \ell_\gamma(\alpha)=\int_X\alpha\wedge\gamma^{n-1}
\]
is strictly positive on its nonzero elements, and the slice of pseudoeffective classes with $\ell_\gamma=1$ is compact. This slice may be empty. Pullback and pushforward by a compact modification are continuous linear maps on Bott--Chern cohomology and preserve pseudoeffectivity.
\end{proposition}
\begin{proof}
Use the finite-dimensional smooth representative space of \cref{lem:constant-trace-bc-representative} to give Bott--Chern cohomology its usual finite-dimensional topology. Any choice of a basis or of the Gauduchon metric gives the same topology. The functional is well-defined by current Stokes and continuous by smooth representatives. If $T\geq0$ represents $\alpha$, then $\ell_\gamma(\alpha)$ is its trace mass. It is nonnegative, and equality forces $T=0$, hence $\alpha=0$. Convexity follows by adding positive currents. If $\alpha$ and $-\alpha$ are both pseudoeffective, their nonnegative masses add to zero; both classes are zero.

Let pseudoeffective classes $\alpha_j$ converge to $\alpha$, and choose positive closed representatives $T_j$. Their masses $\ell_\gamma(\alpha_j)$ are bounded. Weak compactness gives a subsequence converging to a positive closed current $T$. Writing $B_\gamma$ for the continuous smooth representative map of \cref{lem:constant-trace-bc-representative}, continuity of $B_\gamma(T_j)$ shows that the class of $T$ is $\alpha$. Thus the cone is closed. For its mass-one slice, \eqref{eq:constant-trace-representative-bound} bounds $B_\gamma(T)$ in the finite-dimensional representative space. The slice is closed and bounded there, hence compact.

The class maps of a modification are linear and well-defined by the smooth-to-current description of cohomology, so they are continuous between finite-dimensional spaces. Pullback preserves positive closed currents by local psh potentials, as in \cref{prop:closed-current-modification-continuity}; proper pushforward preserves positivity and closedness. These representatives prove the assertions about the cones.
\end{proof}

Compactness of this linear slice does not give a bound for Monge--Amp\`ere masses of potentials. The functional tests one closed current against a fixed smooth form; it is not a top-degree volume of the class.
